% !TeX program = xelatex
\documentclass[12pt]{nextbook}

\UseTemplateSet{
  globalfonts = {libertinus, sanskrit},
  features = {hyperlinks}
}

\newcommand{\ReaderVocabBaselineskip}{15pt}
\newcommand{\ReaderVocabText}[1]{\textcolor{black!35!blue}{#1}}
\newcommand{\ReaderVocabTerm}[1]{#1}
\newcommand{\ReaderVocabSeparator}{}
\usepackage{paracol}
\usepackage{xcolor}
\usepackage{eso-pic}
\usepackage{needspace}

\vbadness=10000
\hbadness=10000
\hfuzz=3pt
\emergencystretch=1.5em

\makeatletter
\AtBeginDocument{%
  \@ifpackageloaded{hyperref}{%
    \pdfstringdefDisableCommands{%
      \def\ruby#1#2{#1}%
      \def\ReaderVocabText#1{#1}%
      \def\ReaderVocabTerm#1{#1}%
    }%
  }{}%
}
\makeatother

\providecommand{\ReaderColumnSep}{6mm}
\providecommand{\ReaderColumnRatio}{0.75}
\providecommand{\ReaderVocabBaselineskip}{20pt}
\providecommand{\ReaderVocabText}[1]{\textbf{#1}}
\providecommand{\ReaderVocabTerm}[1]{\textbf{#1}}
\providecommand{\ReaderVocabSeparator}{:}
\providecommand{\ReaderVocabItem}[2]{\noindent\ReaderVocabTerm{#1}\ReaderVocabSeparator\ #2\par}

\setlength{\columnsep}{\ReaderColumnSep}
\setlength{\columnseprule}{0pt}
\columnratio{\ReaderColumnRatio}
\swapcolumninevenpages

\newcounter{readersection}
\newwrite\vocabfile
\newif\iffirstreaderlabelwritten
\newlength{\readercolumnswidth}
\newlength{\readertitlegutter}
\setlength{\readercolumnswidth}{\dimexpr\textwidth-\columnsep\relax}
\setlength{\readertitlegutter}{.75\columnsep}

\makeatletter
\newcommand{\readerseparatorline}{%
  \if@mainmatter
    \AtTextLowerLeft{%
      \ifodd\value{page}%
        \put(\LenToUnit{\dimexpr .75\readercolumnswidth+.5\columnsep\relax},0){%
          \color{black!35}\rule{0.3pt}{\textheight}%
        }%
      \else
        \put(\LenToUnit{\dimexpr .25\readercolumnswidth+.5\columnsep\relax},0){%
          \color{black!35}\rule{0.3pt}{\textheight}%
        }%
      \fi
    }%
  \fi
}
\makeatother

\AddToShipoutPictureBG{\readerseparatorline}

\newcommand{\currentvocabfile}{\jobname-vocab-\thereadersection.voc}
\newcommand{\beginvocabwrite}{\immediate\openout\vocabfile=\currentvocabfile}
\newcommand{\finishvocabwrite}{\immediate\closeout\vocabfile}
\newcommand{\vocabitem}[2]{\ReaderVocabItem{#1}{#2}}
\newcommand{\printvocablist}{%
  \footnotesize
  \setlength{\parindent}{0pt}%
  \setlength{\parskip}{0pt}%
  \setlength{\baselineskip}{\ReaderVocabBaselineskip}%
  \raggedright
  \input{\currentvocabfile}%
}

\newcommand{\readersectiontitle}[1]{%
  \par
  \begingroup
  \ifodd\value{page}%
    \leftskip=\readertitlegutter%
    \rightskip=0.25\textwidth plus 1fil%
  \else
    \leftskip=\dimexpr .25\readercolumnswidth+.5\columnsep+\readertitlegutter\relax%
    \rightskip=0pt plus 1fil%
  \fi
  \section{#1}%
  \par
  \endgroup
}

\NewDocumentCommand{\voc}{omm}{%
  \ReaderVocabText{#2}%
  \IfNoValueTF{#1}{%
    \immediate\write\vocabfile{\string\vocabitem{#2}{#3}}%
  }{%
    \immediate\write\vocabfile{\string\vocabitem{#1}{#3}}%
  }%
}

\newenvironment{readersection}[1]{%
  \stepcounter{readersection}%
  \beginvocabwrite
  \Needspace{10\baselineskip}%
  \iffirstreaderlabelwritten\else
    \label{reader:firstpage}%
    \global\firstreaderlabelwrittentrue
  \fi
  \readersectiontitle{#1}%
  \begin{paracol}{2}
}{%
  \finishvocabwrite
  \switchcolumn
  \printvocablist
  \end{paracol}
}


\title{संस्कृतपाठावली}
\author{आकाश}
\date{\today}

\begin{document}
\frontmatter
\maketitle
\tableofcontents
\mainmatter
\chapter{१सहस्रपदस्तरः}
\begin{readersection}{नारदस्याश्रमे विद्योपदेशः}
\voc[अथ]{अथ}{then} \voc[नारद]{नारदो}{Nārada} \voc[गिरि]{गिरे}{mountain}\voc[आश्रम]{राश्रम}{hermitage}\voc[उपे]{मुपैति}{to approach}। तत्र \voc[त्रि]{त्रयो}{three} \voc[बाल]{बालाः}{young} पित्रा सह वसन्ति। तेषां \voc[पितृ]{पिता}{father} \voc[वेद]{वेद}{Veda}\voc[विद्]{विदां}{knowing} पुरुषाणां मध्ये वरः। तस्य \voc[भार्या]{भार्या}{wife} \voc[गो]{गाः}{cow} पालयति। आश्रमे \voc[तिल]{तिला}{sesame}\voc[तैल]{स्तैल}{oil}\voc[औषध]{मौषधं}{medicine} च निहितम्। \voc[तु]{तु}{but} बाला नित्यं क्रीडन्ति॥

एकस्मिन्दिने \voc[द्वि]{द्वौ}{two} बालौ \voc[गङ्गा]{गङ्गां}{Ganges} प्रति गच्छतः। \voc[द्वितीय]{द्वितीयो}{second} बालस्तु पितुः समीपे तिष्ठति। तौ गङ्गायां \voc[स्ना]{स्नातु}{to bathe}मिच्छतः। तत्र \voc[सर्प]{सर्प}{snake}\voc[छाया]{श्छायायां}{shadow} शयितः। सर्प एकस्य \voc[पाद]{पादे}{foot} पतति। बालः \voc[पत्]{पतति}{to fall}। तस्य \voc[श्वास]{श्वासो}{breath} मन्दो भवति। \voc[कण्ठ]{कण्ठे}{throat} शब्दो न भवति। \voc[त्वच्]{त्वचि}{skin} \voc[शुक्ल]{शुक्ला}{white} रेखा दृश्यते॥

स बालस्तातेति वदति। पिता नारदं स्मरति। नारदस्तस्य \voc[दृष्टि]{दृष्टिं}{sight} पश्यति। \voc[लोचन]{लोचने}{eye} च पश्यति। ततः स \voc[आचक्ष्]{आचष्टे}{to tell} नायं \voc[कफ]{कफस्य}{phlegm} \voc[हेतु]{हेतुः}{cause} सर्पस्पर्शस्य हेतुरयमिति। पिता शीतलं जलं देहीति वदति। नारदः \voc[शीतल]{शीतल}{cool}मौषधं तैलं चानयेत्याचष्टे॥

भार्या जलमानीय नारदस्य पुरतः स्थापयति। पिता \voc[आसन]{आसनं}{seat} ददाति। नारदस्तत्रोपविशति। स \voc[प्रयोग]{प्रयोगं}{application} \voc[दर्शय्]{दर्शयति}{to show}। \voc[पूर्ववत्]{पूर्वव}{as before}त्पादे तैलं दद्यात्। ततो\voc[अम्ल]{ऽम्ल}{sour}मौषधं \voc[मात्रा]{मात्रया}{measure} देयम्। एतन्नाशनं न सर्पस्य। दुःखस्य शमनं नाम। पूर्णो \voc[नाश]{नाशो}{destruction} न भवति। न च \voc[विनाश]{विनाशः}{destruction} कार्य इति॥

बालस्य श्वासः शनैः समो भवति। पितुश्चेतसि \voc[निश्चय]{निश्चयो}{certainty} जायते। नारदो वदति \voc[धन]{धनं}{wealth} न \voc[मुख्य]{मुख्यं}{chief} विद्यैव धनम्। \voc[विद्या]{विद्या}{knowledge} धना\voc[अधिक]{दधिका}{greater}। \voc[कर्मन्]{कर्म}{action} \voc[शुचि]{शुचि}{pure} चेत्त\voc[वर]{द्वरं}{excellent} भवति। \voc[यतस्]{यतः}{since} शुचिः पुरुषः \voc[पर]{परेषां}{other} हितं करोति \voc[तथा]{तथा}{so} तस्य \voc[शक्ति]{शक्ति}{power}र्वर्धते॥

रात्रौ \voc[याम]{यामे}{watch} \voc[चन्द्र]{चन्द्रो}{moon} गिरिशिखरे दृश्यते। चन्द्रस्य प्रकाश आश्रमं स्पृशति। बालः \voc[स्वप्न]{स्वप्ने}{dream} \voc[विष्णु]{विष्णुं}{Viṣṇu} पश्यामीति वदति। अन्यो बालो \voc[महादेव]{महादेवं}{Mahādeva} पश्यामीति वदति। पिता विष्णुर्महादेवश्च रक्षत इति मन्यते। नारदो वदति देवभक्तिः शुभा किं तु कर्मापि करणीयम्। देवानां स्मरणं चेतसि शुचितां करोति॥

तस्मिन्नाश्रमे सुवर्णाभं किञ्चित्पात्रमस्ति। बालास्त\voc[सुवर्ण]{त्सुवर्णं}{gold} मन्यन्ते। नारदो वदति न सर्वं सुवर्णं य\voc[संकाश]{त्संकाशं}{appearance} सुवर्णवत्। \voc[वज्र]{वज्रं}{thunderbolt} दृढं भवति परं चेतसि \voc[मोह]{मोहो}{delusion} यदि तदपि न रक्षति। मोहः सर्वनाशस्य बीजम्। शुचि चेतः सर्वेषां रक्षणमिति॥

प्रभाते बालः स्वस्थो भवति। नारदो वदति \voc[पञ्चम]{पञ्चमे}{fifth} दिने पुनरौषधं दद्यात्। \voc[द्वादशन्]{द्वादश}{twelve} तिलान्पचेत्। तिलतैलमौषधेन सह योजयेत्। \voc[नित्यम्]{नित्यं}{always} मितं भुञ्जीत। योऽतिमात्रं \voc[भक्षय्]{भक्षयति}{to eat} तस्य \voc[श्रम]{श्रमो}{fatigue} भवति। श्रमात्कफस्य \voc[वृद्धि]{वृद्धि}{increase}र्भवति। कफवृद्ध्या पुनः कण्ठे दोषो भवतीति॥

ततः पिता राज्यस्य चिन्तां करोति। अस्मि\voc[राज्य]{न्राज्ये}{kingdom} \voc[शत्रु]{शत्रवः}{enemy} सन्ति। \voc[पक्ष]{पक्षौ}{side} द्वावपि बलवन्तौ। कः \voc[विषय]{विषयो}{object} मुख्यः कश्च गौण इति पिता पृच्छति। नारदो वदति राज्ञो निश्चयः शुचिर्भवेत्। \voc[उग्र]{उग्रं}{fierce} कर्म राज्यस्य विनाशाय भवति। शत्रुं क्रोधेन न \voc[निहन्]{निहन्यात्}{to kill}। धर्मेणैव जयः कार्य इति॥

पिता पुनः पृच्छति कः \voc[प्रभु]{प्रभुः}{lord} स्यादिति। नारदो वदति यो \voc[मही]{महीं}{earth} मातरं मन्यते स प्रभुः। यो धनं सुवर्णं वज्रं च सर्वस्वं न मन्यते स वरः। प्रजाहितं मुख्यं राज्यकर्म। \voc[पुरा]{पुरा}{formerly} \voc[व्यास]{व्यासो}{Vyāsa} धर्मं कथयामास। \voc[पार्थ]{पार्थो}{Pārtha}ऽपि तच्छुश्राव। तस्माद्राज्ये विद्या धर्मः कर्म च स्थापनीयम्॥

अपराह्णे \voc[वानर]{वानरो}{monkey} वृक्षादवतीर्य आश्रमं प्रविशति। तिलगन्धेन स आगतः। बालाः हसन्ति। भार्या \voc[चित्र]{चित्रो}{variegated} वानर इति वदति। वानरः पादाभ्यां शीघ्रं चलति। नारदो वदति परेषामपि जीवनं प्रियं तस्मादन्नं देयम्। पिता तिला\voc[दा]{न्ददाति}{to give}। वानरस्तिलान्भक्षयित्वा पुनर्गिरिं \voc[व्रज्]{व्रजति}{to go}॥

सायंकाले गङ्गाकथा कथ्यते। पिता वदति गङ्गा \voc[सागर]{सागरं}{ocean} व्रजति। गिरेर्जलं महीं प्राप्नोति। \voc[चन्द्रसंकाश]{चन्द्रसंकाशा}{moonlike} गङ्गा सर्वेषां मनः शीतलं करोति। बालो वदति चन्द्रः सुवर्णसंकाशो दृश्यते। नारदो वदति दृश्यं न सर्वदा सत्यं। दृष्टिर्विषयेषु पतति। विषयास्तु चेतसि मोहं जनयन्ति। शुचिदृष्ट्या विषयाः शुचयो भवन्तीति॥

तस्मिन्काले वैद्य आगच्छति। स \voc[स्थान]{स्थानं}{place} पश्यति। पादस्य स्थानं त्वचं कण्ठं लोचने च पश्यति। सर्पघ्नमौषधं पर्याप्तमिति वदति। नारदो वदति \voc[घ्न]{सर्पघ्नं}{killing} नाम सर्पहन्तृ न स्यात्। पीडाघ्नमत्र मुख्यं भवति। सर्पोऽपि जीवति। न हन्तव्यः। अय\voc[आगम]{मागमो}{tradition} न केवलं वैद्यकस्य धर्मस्यापीति॥

रात्रौ बालाः परस्परं कथयन्ति। \voc[परस्पर]{परस्परं}{mutual} न हास्यं कार्यमिति ज्येष्ठो बालो वदति। यदि कश्चित्पतति तर्हि तं धारयेत्। यदि कस्यचिच्छ्वासो गुरुः स्यात्तर्हि तं पितरमाचक्षीत। यदि \voc[तात]{तातो}{father} न स्यात्तर्हि गुरुमुपेयात्। इत्युक्त्वा ते स्वपन्ति। तेषां स्वप्ने नारदो विष्णुर्महादेवश्च दृश्यन्ते॥

अन्यस्मिन्दिने पिता नारदं पृच्छति। कथं बालानां शिक्षा कार्येति। नारदो वदति वेदविद्या कर्मविद्या चोभे अपेक्षिते। वेदेन चेतः शुचि भवति। कर्मणा लोकहितं भवति। धनदः पुरुषो धनं ददाति। \voc[द]{धनदो}{giving}ऽपि यदि विद्यां न ददाति तर्हि न पूर्णः। विद्यादः पुरुषः परं धनं ददातीति॥

ततः नारदः सर्वानुपदिशति। \voc[अन्तर्]{अन्त}{within}श्चेतसि मोहो न कार्यः। बहिः सुवर्णं दृश्यते चेत्किं तेन। अन्तः शुचिता चेत्सर्वं साधु। \voc[अन्त]{अन्ते}{end} कर्मैव तिष्ठति। \voc[कृत्]{कृतं}{done} शुभं कर्म सुखाय भवति। दुष्कृतं कर्म नाशाय भवति। तस्माच्छुचि कर्म नित्यं कुरुत। एष आगमस्य सार इति॥

प्रातर्नारदो गन्तुमुद्यतः। बालाः तातेति पितरं वदन्ति नारदो मा व्रजत्विति। पिता वदति स पराश्रमं व्रजति। तत्रापि बाला रोगिणो भवन्ति। तत्रापि पितरश्चिन्तयन्ति। तत्रापि भार्या गाः पालयति। तत्रापि तिलास्तैलमौषधं च अपेक्षितम्। नारदः सर्वेषां हिताय व्रजतीति॥

नारदो वदति मम गमनं न विरहाय। विद्याया विस्ताराय। यत्र गच्छामि तत्र वेदः कर्म च कथ्यते। यत्र बालाः सन्ति तत्र मम हृदयं तिष्ठति। यत्र पिता भार्या गाश्च सन्ति तत्र धर्मो जीवति। एवं वदन्नारदो गिरेः पन्थानं व्रजति। बालाः प्रणमन्ति। पिता तिलान्ददाति। भार्या शीतलं जलं ददाति। नारदस्तथैव गच्छति॥
\end{readersection}
\end{document}
